% GENERATED FILE. Edit canonical lessons, then run npm run generate:books.
% canonical-source-hash: fnv1a64:667f88d016c96382
% canonical-lessons: RU-C44-fruit, RU-C44-apple, RU-C44-pear, RU-C44-cherry, RU-C44-berry

\chapter{Fruit, Apple, Pear, Cherry, Berry}
\label{ch:ru-fruit-apple-pear-cherry-berry}
\hlchaptermodality{44}

\begin{chapteropening}
\textbf{By the end of this chapter:} \emph{I can name fruit, an apple, a pear, a cherry and a berry.}

Complete the last lesson of chapter 44: I can name fruit, an apple, a pear, a cherry and a berry.
\end{chapteropening}

\section[\texorpdfstring{\ru{фрукты} (frúkty)}{фрукты (frúkty)}]{\ru{фрукты} (frúkty) — fruit}

\label{lesson:RU-C44-fruit}



\begin{quote}
Before the new one: say the Russian for butter; oil, then the Russian for salt.
\end{quote}

\subsection*{You'll want to know: \ru{фрукты}}
\textbf{\ru{фрукты}} (\emph{frúkty}) — \textquotedblleft{}fruit\textquotedblright{}.

\begin{cousinweb}
\textbf{\ru{фрукты}} came through Polish \emph{frukt} from Latin \emph{frūctus}, \textquotedblleft{}fruit, enjoyment\textquotedblright{}; it is recorded from 1705, in Peter the Great's time (Vasmer). English \emph{fruit} comes from the same Latin word.
\end{cousinweb}

\begin{grammarlens}[title={What you've built}]
One more word for this chapter.
\end{grammarlens}

\subsection*{Your turn}
\begin{itemize}
\raggedright
  \item \emph{Say it:} \emph{frúkty}
  \item \emph{Say it:} \emph{frúkty}, once more
  \item \emph{Say it:} say \emph{frúkty}
  \item \emph{Recall:} say the Russian for butter; oil, then the Russian for salt, then say \emph{frúkty} again
\end{itemize}

\begin{tcolorbox}[breakable,colback=teal!4,colframe=teal!35!black,title={Before you move on}]
What does \ru{фрукты} mean? (\textquotedblleft{}Fruit\textquotedblright{}.) Say it once more.
\end{tcolorbox}
\section[\texorpdfstring{\ru{яблоко} (yábloko)}{яблоко (yábloko)}]{\ru{яблоко} (yábloko) — an apple}

\label{lesson:RU-C44-apple}



\begin{quote}
Before the new one: say the Russian for salt, then the Russian for fruit.
\end{quote}

\subsection*{You'll want to know: \ru{яблоко}}
\textbf{\ru{яблоко}} (\emph{yábloko}) — \textquotedblleft{}an apple\textquotedblright{}.

\begin{cousinweb}
\textbf{\ru{яблоко}} $\leftarrow$ Proto-Slavic \emph{*ablŭko} — inherited. Vasmer lists Lithuanian \emph{obuolỹs}, Old High German \emph{apful} and Old Irish \emph{aball}, all \textquotedblleft{}apple\textquotedblright{}: English \emph{apple} is a cousin.
\end{cousinweb}

\begin{grammarlens}[title={What you've built}]
One more word for this chapter.
\end{grammarlens}

\subsection*{Your turn}
\begin{itemize}
\raggedright
  \item \emph{Say it:} \emph{yábloko}
  \item \emph{Say it:} \emph{yábloko}, once more
  \item \emph{Say it:} say \emph{yábloko}
  \item \emph{Recall:} say the Russian for salt, then the Russian for fruit, then say \emph{yábloko} again
\end{itemize}

\begin{tcolorbox}[breakable,colback=teal!4,colframe=teal!35!black,title={Before you move on}]
What does \ru{яблоко} mean? (\textquotedblleft{}An apple\textquotedblright{}.) Say it once more.
\end{tcolorbox}
\section[\texorpdfstring{\ru{груша} (grúsha)}{груша (grúsha)}]{\ru{груша} (grúsha) — a pear}

\label{lesson:RU-C44-pear}



\begin{quote}
Before the new one: say the Russian for fruit, then the Russian for an apple.
\end{quote}

\subsection*{You'll want to know: \ru{груша}}
\textbf{\ru{груша}} (\emph{grúsha}) — \textquotedblleft{}a pear\textquotedblright{}.

\begin{cousinweb}
\textbf{\ru{груша}} — probably a very old loan: the first sound wavers between \emph{\ru{г}-} and \emph{\ru{к}-} (Bulgarian \emph{\ru{круша}}), and Lithuanian \emph{kriáušė}, \textquotedblleft{}pear\textquotedblright{}, has the same source (Vasmer). Vasmer reports a comparison with a Kurdish word, but the donor is unknown.
\end{cousinweb}

\begin{grammarlens}[title={What you've built}]
One more word for this chapter.
\end{grammarlens}

\subsection*{Your turn}
\begin{itemize}
\raggedright
  \item \emph{Say it:} \emph{grúsha}
  \item \emph{Say it:} \emph{grúsha}, once more
  \item \emph{Say it:} say \emph{grúsha}
  \item \emph{Recall:} say the Russian for fruit, then the Russian for an apple, then say \emph{grúsha} again
\end{itemize}

\begin{tcolorbox}[breakable,colback=teal!4,colframe=teal!35!black,title={Before you move on}]
What does \ru{груша} mean? (\textquotedblleft{}A pear\textquotedblright{}.) Say it once more.
\end{tcolorbox}
\section[\texorpdfstring{\ru{вишня} (víshnya)}{вишня (víshnya)}]{\ru{вишня} (víshnya) — a cherry}

\label{lesson:RU-C44-cherry}



\begin{quote}
Before the new one: say the Russian for an apple, then the Russian for a pear.
\end{quote}

\subsection*{You'll want to know: \ru{вишня}}
\textbf{\ru{вишня}} (\emph{víshnya}) — \textquotedblleft{}a cherry\textquotedblright{}.

\begin{cousinweb}
\textbf{\ru{вишня}} $\leftarrow$ Proto-Slavic \emph{*višĭnja} — inherited. Vasmer relates it to Old High German \emph{wîhsila}, \textquotedblleft{}sour cherry\textquotedblright{}, and through it to Latin \emph{viscum}, \textquotedblleft{}birdlime\textquotedblright{}, because cherry gum and mistletoe berries were both used to snare birds.
\end{cousinweb}

\begin{grammarlens}[title={What you've built}]
One more word for this chapter.
\end{grammarlens}

\subsection*{Your turn}
\begin{itemize}
\raggedright
  \item \emph{Say it:} \emph{víshnya}
  \item \emph{Say it:} \emph{víshnya}, once more
  \item \emph{Say it:} say \emph{víshnya}
  \item \emph{Recall:} say the Russian for an apple, then the Russian for a pear, then say \emph{víshnya} again
\end{itemize}

\begin{tcolorbox}[breakable,colback=teal!4,colframe=teal!35!black,title={Before you move on}]
What does \ru{вишня} mean? (\textquotedblleft{}A cherry\textquotedblright{}.) Say it once more.
\end{tcolorbox}
\section[\texorpdfstring{\ru{ягода} (yágoda)}{ягода (yágoda)}]{\ru{ягода} (yágoda) — a berry}

\label{lesson:RU-C44-berry}



\begin{quote}
Before the new one: say the Russian for a pear, then the Russian for a cherry.
\end{quote}

\subsection*{You'll want to know: \ru{ягода}}
\textbf{\ru{ягода}} (\emph{yágoda}) — \textquotedblleft{}a berry\textquotedblright{}.

\begin{cousinweb}
\textbf{\ru{ягода}} $\leftarrow$ Proto-Slavic \emph{*agoda}, \emph{*aga} with an added \emph{-oda} — inherited, a cousin of Lithuanian \emph{úoga}, \textquotedblleft{}berry\textquotedblright{} (Vasmer). Russian added \emph{j-} at the front, as in \emph{\ru{яблоко}}, \textquotedblleft{}apple\textquotedblright{}.
\end{cousinweb}

\begin{grammarlens}[title={What you've built}]
One more word for this chapter.
\end{grammarlens}

\subsection*{Your turn}
\begin{itemize}
\raggedright
  \item \emph{Say it:} \emph{yágoda}
  \item \emph{Say it:} \emph{yágoda}, once more
  \item \emph{Say it:} say \emph{yágoda}
  \item \emph{Recall:} say the Russian for a pear, then the Russian for a cherry, then say \emph{yágoda} again
\end{itemize}

\begin{tcolorbox}[breakable,colback=teal!4,colframe=teal!35!black,title={Before you move on}]
What does \ru{ягода} mean? (\textquotedblleft{}A berry\textquotedblright{}.) Say it once more.
\end{tcolorbox}
